Efficiently learning goal-conditioned policies for robotic manipulation remains a significant challenge, particularly in leveraging high-dimensional visual data and natural language instructions. Furthermore, reward shaping has been a long-standing issue in reinforcement learning. This thesis introduces a novel framework that learns goal-conditioned value functions using a self-supervised energy-based model (EBM). The approach is founded on a Joint-Embedding Predictive Architecture (JEPA), which learns to predict future states in a latent representation space.

The model's context encoder processes a history of the last k video frames to form a rich representation of the current state. Conditioned on this context and a latent goal vector derived from text instructions, a transformer-based predictor module anticipates the latent patch encodings of the next frame. The model is trained entirely in a self-supervised manner, minimizing the distance between its prediction and the actual future frame's representation, which is generated by a target encoder, whose weights are updated as the EMA of the context encoder.

The prediction loss itself can be interpreted as an ``energy'' value at inference time. This energy, representing the discrepancy between the model's expectation and the observed reality, serves as an implicit, goal-conditioned value function. A low-energy state (i.e., low prediction error) signifies that the observed transition is highly probable given the history and the language command, effectively signalling progress toward the goal. This JEPA framework enables the learning of complex manipulation behaviors from raw visual data, improves sample efficiency by removing the need for explicit rewards, and provides a flexible mechanism for constructing dense reward models for goal-conditioned reinforcement learning.

The black-box nature of such neural architectures has long been a hurdle to the widespread adoption of generalist robotic systems powered by large foundational models, such as vision-language models (VLAs). This calls for a mechanism to make these models more explainable and interpretable, thus making them robust and trustworthy. Our goal-conditioned energy-based value function, being a differentiable module, allows the inspection of gradients through the different nodes of the computation graph. Analyses of gradients and higher-order differentials of different patches or pixels in the image can provide us insights into the implicit reasoning of the energy-based model, thus enabling inspection of the contribution of different parts of the input to the discrepancy, \textit{aka} energy.
